\documentclass[12pt, twoside]{article}
\input{../preamble}
\usepackage{float}

\title{IB}
\author{Chris Huson}
\date{March 2026}

\fancyhead[LE]{\thepage}
\fancyhead[LO]{La Scuola d'Italia / Huson / IB Math: Series \\* 20 March 2026}

\fancypagestyle{firstpage}{%
  \fancyhf{}%
  \fancyhead[RO]{First \& last name: \hspace{2.25cm} \,\\ \,}%
  \fancyhead[LO]{La Scuola d'Italia / Huson / IB Math: Series \\* 20 March 2026}%
}

\begin{document}
\thispagestyle{firstpage}

\subsubsection*{1.3 Classwork: Series and Sigma Notation}
\begin{enumerate}[itemsep=0.5cm]
\item An arithmetic series has first term $u_1=3$ and common difference $d=5$.
    \begin{enumerate}[itemsep=0.25cm]
        \item Find $u_{20}$, the twentieth term. \hfill [2]
        \item Find $S_{20}$, the sum of the first 20 terms. \hfill [2]
        \item The $n^{th}$ term of the sequence is 148. Find $n$. \hfill [2]
    \end{enumerate}
    \begin{tikzpicture}
        \draw (0,1) rectangle (15.5,11);
        \draw [dotted] (1,10)--(14,10);
        \draw [dotted] (1,9)--(14,9);
        \draw [dotted] (1,8)--(14,8);
        \draw [dotted] (1,7)--(14,7);
        \draw [dotted] (1,6)--(14,6);
        \draw [dotted] (1,5)--(14,5);
        \draw [dotted] (1,4)--(14,4);
        \draw [dotted] (1,3)--(14,3);
    \end{tikzpicture}

\newpage
\item A geometric sequence has first term $u_1=6$ and common ratio $\displaystyle r=\frac{1}{2}$.
    \begin{enumerate}[itemsep=0.25cm]
        \item Write down the first four terms of the sequence. \hfill [2]
        \item Find $S_8$, the sum of the first 8 terms. \hfill [2]
        \item Explain why the infinite sum $S_\infty$ exists, and find its value. \hfill [3]
    \end{enumerate}
    \begin{tikzpicture}
        \draw (0,1) rectangle (15.5,11);
        \draw [dotted] (1,10)--(14,10);
        \draw [dotted] (1,9)--(14,9);
        \draw [dotted] (1,8)--(14,8);
        \draw [dotted] (1,7)--(14,7);
        \draw [dotted] (1,6)--(14,6);
        \draw [dotted] (1,5)--(14,5);
        \draw [dotted] (1,4)--(14,4);
        \draw [dotted] (1,3)--(14,3);
    \end{tikzpicture}

\newpage
\item Consider the arithmetic sequence $7, 11, 15, 19, \ldots$
    \begin{enumerate}[itemsep=0.25cm]
        \item Find the general term $u_n$. \hfill [2]
        \item Find $\displaystyle \sum_{k=1}^{15} u_k$. \hfill [2]
        \item Find the least value of $n$ such that $S_n > 1000$. \hfill [3]
    \end{enumerate}
    \begin{tikzpicture}
        \draw (0,1) rectangle (15.5,11);
        \draw [dotted] (1,10)--(14,10);
        \draw [dotted] (1,9)--(14,9);
        \draw [dotted] (1,8)--(14,8);
        \draw [dotted] (1,7)--(14,7);
        \draw [dotted] (1,6)--(14,6);
        \draw [dotted] (1,5)--(14,5);
        \draw [dotted] (1,4)--(14,4);
        \draw [dotted] (1,3)--(14,3);
    \end{tikzpicture}

\newpage
\item The sum of the first $n$ terms of a geometric series is given by
\[S_n = 4\left(1 - \left(\frac{1}{3}\right)^n\right)\]
    \begin{enumerate}[itemsep=0.25cm]
        \item Find $S_1$ and hence write down $u_1$. \hfill [2]
        \item Find $S_2$ and hence find $u_2$. \hfill [2]
        \item Find the common ratio $r$. \hfill [1]
        \item Find $S_\infty$. \hfill [2]
    \end{enumerate}
    \begin{tikzpicture}
        \draw (0,1) rectangle (15.5,10);
        \draw [dotted] (1,9)--(14,9);
        \draw [dotted] (1,8)--(14,8);
        \draw [dotted] (1,7)--(14,7);
        \draw [dotted] (1,6)--(14,6);
        \draw [dotted] (1,5)--(14,5);
        \draw [dotted] (1,4)--(14,4);
        \draw [dotted] (1,3)--(14,3);
    \end{tikzpicture}

\end{enumerate}
\end{document}
